% !PRINTED 26
\weberpage{26}
%<<CONT
and the functions $\Phi$, $\Psi$ take the canonical form:
\weq{14}{\begin{aligned}
\Phi &= \Sigma\,\beta_1^2\alpha_2\alpha_3\alpha_4\,y_1^2,\\
\Psi &= \Sigma\,\alpha_1^2\beta_2\beta_3\beta_4\,y_1^2.
\end{aligned}}
The determinant of the system (12), (14), regarded as linear equations
for $y_1^2$, $y_2^2$, $y_3^2$, $y_4^2$, is
\weq{15}{\begin{vmatrix}
\alpha_1^2\beta_2\beta_3\beta_4, & \alpha_2^2\beta_1\beta_3\beta_4, & \alpha_3^2\beta_1\beta_2\beta_4, & \alpha_4^2\beta_1\beta_2\beta_3\\
\beta_1^2\alpha_2\alpha_3\alpha_4, & \beta_2^2\alpha_1\alpha_3\alpha_4, & \beta_3^2\alpha_1\alpha_2\alpha_4, & \beta_4^2\alpha_1\alpha_2\alpha_3\\
\alpha_1, & \alpha_2, & \alpha_3, & \alpha_4\\
\beta_1, & \beta_2, & \beta_3, & \beta_4
\end{vmatrix}.}
It vanishes when $\alpha_1:\alpha_2 = \beta_1:\beta_2$, because the
first two columns then become proportional to one another, and likewise
when $\alpha_i:\alpha_k = \beta_i:\beta_k$. Therefore (15), regarded as
an integral function of the $\alpha$, $\beta$, is divisible by
\weq{16}{\alpha_i\beta_k - \alpha_k\beta_i = (\alpha_i\beta_k),}
and consequently also by the product of all these factors
\weq{17}{(\alpha_1\beta_2)(\alpha_1\beta_3)(\alpha_1\beta_4)(\alpha_2\beta_3)(\alpha_2\beta_4)(\alpha_3\beta_4) = \Delta.}
From the comparison of the degrees it follows that the two functions
(15), (17) can differ only by a numerical factor, and this factor is
found, from the assumption $\beta_1 = 0$, $\alpha_2 = 0$,
$\alpha_4 = 0$, to be equal to 1. The determinant (15) is therefore
equal to the product $\Delta$, and the square of $\Delta$ is the
discriminant $D$ of the function $f(\xi,\eta)$, hence equal to the
function $D$, which we have assumed to be different from zero. The
$y_1^2$, $y_2^2$, $y_3^2$, $y_4^2$ can therefore be represented
linearly by $\varphi$, $\psi$, $\Phi$, $\Psi$.

\begin{center}
\rule{2cm}{0.4pt}
\end{center}

The covariants $\Phi$, $\Psi$ change when $\varphi$ and $\psi$ are
replaced according to (4) by $\varphi'$, $\psi'$. They therefore belong
not to the fundamental curve, but to the pair of forms $\varphi$,
$\psi$. But under this substitution $\Phi$ and $\Psi$, as linear
functions of $y_1^2$, $y_2^2$, $y_3^2$, $y_4^2$, pass over into linear
functions of $\varphi$, $\psi$, $\Phi$, $\Psi$. These expressions are
rather complicated. But we need them here only for the points of the
fundamental curve, that is, for $\varphi = 0$, $\psi = 0$, and under
this hypothesis they become very simple. For in that case we have
\weq{18}{\begin{aligned}
\Phi' &= \Sigma\,(p\alpha_1+q\beta_1)^2(m\alpha_2+n\beta_2)(m\alpha_3+n\beta_3)(m\alpha_4+n\beta_4)\,y_1^2,\\
\Psi' &= \Sigma\,(m\alpha_1+n\beta_1)^2(p\alpha_2+q\beta_2)(p\alpha_3+q\beta_3)(p\alpha_4+q\beta_4)\,y_1^2,
\end{aligned}}
%CONT>>
